\subsection{Generalization to complexes of multimodules}

Let B and B' be two rings, C a complex of (B, A)-bimodules, C' a complex of (A, B')-bimodules (X, p. 43); then (C ⊗\_A C', D) (X, p. 61) is a complex of (B, B')-bimodules and the canonical homomorphism

\[
\gamma : \mathrm{H} (\mathbf {C}) \otimes_ {\mathbf {A}} \mathrm{H} (\mathbf {C} ^ {\prime}) \rightarrow \mathrm{H} (\mathbf {C} \otimes_ {\mathbf {A}} \mathbf {C} ^ {\prime})
\]

is compatible with the structures of (B, B')-bimodules of the two members.

If B'' is a third ring, and C'' a complex of (B', B'')-bimodules, the canonical homomorphism (II, p. 64, prop. 8)

\[
\left(\mathrm{C} \otimes_ {\mathrm{A}} \mathrm{C} ^ {\prime}\right) \otimes_ {\mathrm{B} ^ {\prime}} \mathrm{C} ^ {\prime \prime} \rightarrow \mathrm{C} \otimes_ {\mathrm{A}} \left(\mathrm{C} ^ {\prime} \otimes_ {\mathrm{B} ^ {\prime \prime}} \mathrm{C} ^ {\prime \prime}\right)
\]

is an isomorphism of complexes of (B, B'')-bimodules.

More generally, we leave it to the reader to develop the theory of tensor products of finite, totally ordered families of complexes of multimodules, following the model of no. 1 (X, p. 63) and of II, pp. 65 to 72 (associativity and commutativity isomorphisms, \ldots).

Let B and B′ be two rings, M a (B, A)-bimodule, N an (A, B′)-bimodule: then L(M) ⊗\_A L(N) is a complex of (B, B')-bimodules, so that Tor\^{}A (M, N) is endowed with a natural structure of graded (B, B')-bimodule; on the degree 0 term, this structure coincides with that of M ⊗\_A N.

If \(\lambda\in\mathbf{B},\lambda^{\prime}\in\mathbf{B}^{\prime}\) , and if we denote by \(\lambda_{\mathbf{M}},\lambda_{\mathbf{N}}^{\prime},\lambda_{\mathbf{T}},\lambda_{\mathbf{T}}^{\prime}\) , the homotheties \(x\mapsto\lambda x,y\mapsto y\lambda^{\prime},z\mapsto\lambda z,z\mapsto z\lambda^{\prime}\) of \(\mathbf{M},\mathbf{N},\mathrm{Tor}^{\mathbf{A}}(\mathbf{M},\mathbf{N}),\mathrm{Tor}^{\mathbf{A}}(\mathbf{M},\mathbf{N})\) respectively, then

\[
\lambda_ {\mathrm{T}} = \operatorname{Tor} ^ {\mathbf {A}} \left(\lambda_ {\mathbf {M}}, 1 _ {\mathbf {N}}\right), \quad \lambda_ {\mathrm{T}} ^ {\prime} = \operatorname{Tor} ^ {\mathbf {A}} \left(1 _ {\mathbf {M}}, \lambda_ {\mathbf {N}} ^ {\prime}\right),
\]

which provides another description of the bimodule structure of Tor \(^{A}\) (M, N). We leave it to the reader to generalize nos. 5 and 7 to the case of complexes of multimodules.

